\documentclass{book}
\usepackage{graphicx} % Required for inserting images
\usepackage[OT4]{fontenc}
\usepackage[polish]{babel}
\usepackage{geometry}
\usepackage{url}
\usepackage{amsmath}
\usepackage[backend=biber, style=numeric]{biblatex}

\geometry{top=2cm, bottom=2cm, left=2.5cm, right=2cm, headsep=1.25cm, footskip=1.25cm}

\title{Kompresja wag modeli sztucznej inteligencji - metody i wpływ}
\author{Jakub Grula}
\date{February 2026}

\addbibresource{bib.bib}

\begin{document}

\begin{titlepage}
    \centering
    \textbf{\Large Katolicki Uniwersytet Lubelski Jana Pawła II}\\[0.5cm]
    \textbf{\Large Wydział Filozofii}\\[0.5cm]
    \textbf{\Large Instytut Filozofii}\\[2.5cm]
    
    \Large{Sztuczna inteligencja, studia stacjonarne I stopnia}\\[2cm]
    
    \textbf{\Large Jakub Grula}\\[0.5cm]
    \Large{Nr albumu: 162321}\\[2.5cm]
    
    \textbf{\LARGE Kompresja wag modeli sztucznej inteligencji\\ - metody i wpływ}\\

    \vfill
    
    \begin{flushright}
    {
        \Large{Praca licencjacka}\\[0.1cm]
        \Large{napisana na seminarium} 
        \textbf{\Large Metody i zastosowania uczenia maszynowego}\\[0.1cm]
        \Large{pod kierunkiem}\\[0.1cm]
        \textbf{\Large dr. hab. inż. Krzysztofa Pancerza, prof. KUL}\\[3.5cm]
        }
    \end{flushright}

    
    \textbf{\Large Lublin 2026}\\[1cm]
\end{titlepage}

\tableofcontents
\chapter{Wstęp}
Rozwój modeli sztucznej inteligencji pędzi bez opamiętania. W ciągu kilku ostatnich lat wszyscy byliśmy świadkami niesamowitych postępów w dziedzinach uczenia maszynowego, głębokiego uczenia i przetwarzania języka naturalnego. Coraz więcej osób przekonuje się do różnego rodzaju modeli SI i używa ich w coraz szerszych obszarach swojego życia. Niestety wraz z tym wszystkim dobre modele niespecjalnie stają się mniejsze, wręcz przeciwnie. Wiele z nich jest ogromnych rozmiarów przez ciągle zwiększające się ilości parametrów i nieustannie komplikujące się architektury, co utrudnia ich przechowywanie, przenoszenie oraz wczytywanie do pamięci urządzenia. Dla przykładu, najnowszy state-of-the-art model z rodziny Qwen zajmuje kilkaset gigabajtów \cite{qwen2025technical}. Wspomniany problem dotyczy nie tylko przestrzeni pamięci, czy to dyskowej, czy sprzętowej, ale też ma wpływ na wydajność, szybkość, a także na koszty uruchomienia modelu. W odpowiedzi na ten problem inżynierowie pracują nad różnymi metodami kompresji modeli, które pozwalają na zmniejszenie ich rozmiaru o różnym stopniu straty jakości. Jedną z takich metod jest kwantyzacja, czyli proces, w którym wartości parametrów modelu są reprezentowane za pomocą mniejszej liczby bitów, co pozwala na znaczne zmniejszenie rozmiaru modelu. Kwantyzacja może pomagać zmniejszyć zajmowane miejsce na dysku przez model, ale jest specyficzną formą kompresji, która niekoniecznie wymaga kroku dekompresji. Oznacza to, że modele skompresowane za pomocą niektórych kwantyzacji mogą być bezpośrednio używane bez potrzeby wcześniejszego dekompresowania, co jest dużą zaletą w kontekście wydajności i szybkości działania. \\

W kolejnych rozdziałach pracy przejdę przez ważne opisy, terminy i definicje, które pozwolą zrozumieć dokładniej, czym jest kwantyzacja, jakie dokładnie są jej rodzaje, korzyści oraz koszty. Następnie przedstawiony zostanie mój autorski przykład implementacji kwantyzacji wraz z wynikami testów na różnego rodzaju modelach w różnych konfiguracjach kompresji. Ostatnia część jest poświęcona publicznym rozwiązaniom i formatom, z~których korzystają najbardziej popularne modele i korporacje.
\chapter{Definicje i podstawy}
\section{Definicje i podstawy}

Zanim zostaną omówione konkretne metody kompresji wag modeli sztucznej inteligencji, warto uporządkować podstawowe pojęcia związane z reprezentacją danych numerycznych oraz samą kwantyzacją. W literaturze termin ten pojawia się w wielu odmianach, jednak w kontekście modeli uczenia maszynowego oznacza on przede wszystkim redukcję precyzji zapisu parametrów tak, aby zmniejszyć zużycie pamięci, poprawić przepustowość pamięciową i skrócić czas wykonywania obliczeń \cite{jacob2017quantization,dettmers2023case4bit}.

\subsection{Tensory, parametry i warstwy modelu}

Współczesne modele sieci neuronowych opisuje się najczęściej za pomocą tensorów, czyli wielowymiarowych tablic liczb. W praktyce najważniejsze są dwa rodzaje takich danych: aktywacje, które są wyliczane w trakcie działania modelu, oraz parametry, czyli wagi i biasy zapisane po zakończeniu treningu. To właśnie parametry zajmują znaczną część pamięci modelu i stanowią główny obszar zainteresowania przy jego kompresji.

Model składa się z warstw, a każda warstwa realizuje określoną operację matematyczną. W warstwach gęstych, konwolucyjnych i rekurencyjnych dominują mnożenia macierzy i sumowania wartości, dlatego nawet niewielka zmiana sposobu reprezentacji wag może istotnie wpłynąć zarówno na rozmiar modelu, jak i na jego wydajność obliczeniową. W dalszej części pracy pojęcia \textit{waga}, \textit{parametr} i \textit{wartość wagowa} są używane zamiennie, o ile z~kontekstu nie wynika inaczej.

\subsection{Reprezentacja liczb}

W pamięci komputera liczby mogą być zapisywane w różnych formatach. Najczęściej spotykaną reprezentacją w modelach uczenia maszynowego jest format zmiennoprzecinkowy, np. FP32, FP16 lub BF16. Format FP32 zapewnia dużą precyzję, ale zajmuje 32 bity na jedną liczbę. Z kolei formaty o mniejszej liczbie bitów, takie jak FP16, zmniejszają rozmiar danych, lecz kosztem dokładności i zakresu reprezentowanych wartości.

W przypadku kwantyzacji dąży się zazwyczaj do przejścia z reprezentacji zmiennoprzecinkowej do reprezentacji o niższej precyzji, najczęściej całkowitej. Przykładowo, liczby mogą być przechowywane w postaci INT8 albo INT4. Taki zapis jest znacznie bardziej ekonomiczny pamięciowo, lecz wymaga dodatkowego mechanizmu odwzorowującego wartości ciągłe na dyskretne poziomy kwantyzacji. Mechanizm ten opiera się zwykle na współczynniku skali oraz, w niektórych wariantach, na punkcie zerowym.

\subsection{Formaty niskiej precyzji}

W praktyce współczesnych modeli najczęściej spotyka się trzy grupy formatów niskiej precyzji. Pierwszą są klasyczne liczby całkowite, takie jak INT8 i INT4, które stanowią podstawę wielu metod kwantyzacji po treningu \cite{jacob2017quantization,dettmers2023case4bit}.

Drugą grupę stanowią formaty zmiennoprzecinkowe o ograniczonej liczbie bitów, takie jak FP16 i BF16. Wykorzystuje się je tam, gdzie potrzeba kompromisu między dokładnością a rozmiarem danych. W praktyce BF16 jest szczególnie ważny w nowoczesnym treningu i wnioskowaniu modeli, ponieważ zachowuje szeroki zakres wartości przy niższych kosztach pamięciowych.

Trzecią grupę stanowią formaty wyspecjalizowane, takie jak NF4, opisany w pracach o QLoRA, oraz techniki hybrydowe dla 8-bitowego wnioskowania, spopularyzowane przez LLM.int8() \cite{dettmers2023qlora,dettmers2022llmint8}.

\subsection{Uczenie, wnioskowanie i dostrajanie}

Uczenie modelu polega na takim dopasowywaniu jego parametrów, aby minimalizować błąd względem danych treningowych. W praktyce oznacza to wielokrotne przetwarzanie przykładów wejściowych, obliczanie funkcji straty i aktualizację wag. Wnioskowanie jest natomiast etapem użycia już wytrenowanego modelu do generowania odpowiedzi, prognoz lub klasyfikacji.

W pracy pojawia się również pojęcie dostrajania, czyli fine-tuningu. Oznacza ono dalsze uczenie modelu, który został już wcześniej wstępnie wytrenowany, zwykle na mniejszym i bardziej wyspecjalizowanym zbiorze danych. Taki model przejmuje ogólne własności nabyte wcześniej, ale adaptuje się do nowego zadania.

W kontekście metod kompresji ważne jest też pojęcie kalibracji. Jest to wykorzystanie reprezentatywnego, zwykle niewielkiego zbioru danych do oszacowania parametrów potrzebnych do kwantyzacji, takich jak skala, rozkład aktywacji lub znaczenie poszczególnych kanałów. Kalibracja nie jest pełnym treningiem, ale może znacząco wpływać na końcową jakość skompresowanego modelu.

W literaturze kalibracja często pojawia się razem z podejściem \textit{post-training quantization}, gdzie model nie jest ponownie uczony, lecz tylko dostrajany numerycznie na podstawie ograniczonej próbki danych \cite{frantar2023gptq,NIPS1989_6c9882bb}.

\subsection{Rodzaje zadań uczenia maszynowego}

W pracy pojawiają się trzy podstawowe typy zadań: regresja, klasyfikacja binarna i klasyfikacja wieloklasowa. Regresja polega na przewidywaniu wartości liczbowej, na przykład ceny, ryzyka lub innego parametru ciągłego. Klasyfikacja binarna wybiera jedną z dwóch klas, na przykład tak lub nie. Klasyfikacja wieloklasowa przypisuje przykład do jednej z wielu klas, jak w przypadku rozpoznawania słowa ze ścieżki audio albo klasy przynależności obrazu.

Takie rozróżnienie ma znaczenie nie tylko dla doboru metryki oceny, ale również dla samej konstrukcji modelu. Inne funkcje aktywacji, inne warstwy wyjściowe i inne funkcje straty są naturalne dla regresji, a inne dla klasyfikacji.

\subsection{Warstwy i funkcje aktywacji}

Warstwa Dense, zwana też warstwą gęstą, łączy każdy element wejścia z każdym elementem wyjścia za pomocą macierzy wag. Jest to jedna z najprostszych i najczęściej stosowanych warstw w modelach klasyfikacyjnych oraz regresyjnych.

Warstwy Conv1D i Conv2D realizują operację konwolucji, czyli przesuwania filtra po sygnale jednowymiarowym lub dwuwymiarowym (np. obrazie). Są szczególnie przydatne w analizie danych sekwencyjnych, audio i~obrazów, ponieważ potrafią wykrywać lokalne wzorce.

Warstwa GRU należy do rodziny sieci rekurencyjnych i służy do przetwarzania danych sekwencyjnych, w~których istotna jest kolejność elementów. W przeciwieństwie do prostych warstw gęstych uwzględnia ona stan pamięci z poprzednich kroków.

LayerNormalization to warstwa normalizacyjna, która przekształca wartości wewnątrz próbki tak, aby poprawić stabilność uczenia i ułatwić propagację gradientu. Tego typu warstwy często pojawiają się w głębokich architekturach jako sposób na ustabilizowanie przebiegu treningu.

Funkcja aktywacji określa, jak neurony przekształcają otrzymane sygnały. W pracy pojawiają się między innymi GeLU, Sigmoid i Softmax. GeLU jest gładką funkcją aktywacji często stosowaną w nowoczesnych architekturach. Sigmoid odwzorowuje wartości do przedziału od 0 do 1 i dlatego nadaje się do klasyfikacji binarnej. Softmax przekształca wektor wyjściowy w rozkład prawdopodobieństwa po klasach i jest typową funkcją wyjściową w klasyfikacji wieloklasowej.

GeLU zostało opisane jako \textit{Gaussian Error Linear Unit} i pojawia się w licznych nowoczesnych modelach głębokiego uczenia, dlatego w tej pracy traktuję je jako standardową funkcję aktywacji dla głębokich warstw gęstych \cite{hendrycks2023gaussian}.

\subsection{Metryki oceny modeli}

Do porównywania modeli wykorzystuje się metryki, czyli miary jakości predykcji. Dokładność określa, jaki odsetek wszystkich odpowiedzi został przewidziany poprawnie. Precyzja mówi, jaka część przykładów oznaczonych przez model jako pozytywne rzeczywiście taka była. Czułość informuje, jaka część wszystkich rzeczywiście pozytywnych przykładów została wykryta. Miara F1 łączy precyzję i czułość w jedną miarę harmoniczną, dzięki czemu dobrze nadaje się do oceny modeli przy niezbalansowanych danych.

W zadaniach regresyjnych stosuje się inne miary. MAE, czyli średni błąd bezwzględny, pokazuje przeciętną wartość odchylenia predykcji od wartości rzeczywistej. MSE, czyli średni błąd kwadratowy, silniej karze większe odchylenia. RMSE jest pierwiastkiem z MSE i ma tę samą jednostkę co wartość przewidywana, dzięki czemu łatwiej go interpretować.

\subsection{Pojęcia obliczeniowe}

W rozdziale o implementacji pojawiają się także pojęcia związane z realizacją obliczeń. CPU oznacza procesor centralny, GPU jest procesorem graficznym przystosowanym do równoległego wykonywania wielu prostych operacji, a CUDA jest platformą i zestawem narzędzi umożliwiających wykonywanie obliczeń ogólnego przeznaczenia na kartach NVIDIA. Kernel CUDA to pojedyncza funkcja uruchamiana równolegle przez wiele wątków GPU.

CUDA w tej pracy odnosi się do modelu programowania GPU i do zestawu narzędzi opisanych przez NVIDIA, który pozwala kompilować oraz uruchamiać kod wykonywany na kartach graficznych \cite{nvidia2024cuda}.

Pojęcie \textit{batch size} oznacza liczbę przykładów przetwarzanych w jednym kroku obliczeń. W większych \textit{batchach} (partiach) część operacji da się lepiej zrównoleglić, ale rośnie też zapotrzebowanie na pamięć. W przypadku tej pracy istotny jest również rozmiar bloku, czyli liczba wag grupowanych razem podczas kwantyzacji. Ten parametr wpływa bezpośrednio na stosunek jakości modelu do rozmiaru pliku wynikowego.

\subsection{Istota kwantyzacji}

Najprościej rzecz ujmując, kwantyzacja polega na przypisaniu liczb z przestrzeni rzeczywistej do skończonego zbioru wartości reprezentowalnych w wybranym formacie. Jeżeli oznaczymy oryginalną wartość przez $x$, to proces kwantyzacji można opisać funkcją
\begin{equation}
	q = Q(x),
\end{equation}
gdzie $q$ jest wartością skwantyzowaną, a $Q$ to funkcja kwantyzacji. Operacja odwrotna, czyli dekwantyzacja, przybliża wartość oryginalną na podstawie wartości skwantyzowanej:
\begin{equation}
	\hat{x} = D(q),
\end{equation}
gdzie $\hat{x}$ jest przybliżeniem oryginalnej wartości $x$, a $D$ to funkcja dekwantyzacji.

W prostym wariancie symetrycznym przyjmuje się, że skwantyzowana wartość powstaje przez podzielenie liczby przez skalę $s$, zaokrąglenie i ewentualne ograniczenie zakresu do dozwolonych granic. Dla formatu o $b$~bitach zakres ten można zapisać jako
\begin{equation}
	\left[-2^{b-1},\, 2^{b-1}-1\right].
\end{equation}
W praktyce spotyka się także warianty, w których zakres jest przesunięty lub kodowany w sposób niestandardowy, jednak idea pozostaje taka sama: odwzorować wartości rzeczywiste na mniejszą liczbę poziomów dyskretnych.

\subsection{Skala i punkt zerowy}

W kwantyzacji symetrycznej najważniejszym parametrem jest skala. Umożliwia ona przeliczenie wartości z domeny oryginalnej na domenę skwantyzowaną i z powrotem. Jeśli blok wag ma wartości o największej wartości bezwzględnej równej $\max(|x|)$, to klasyczny dobór skali dla $b$ bitów można zapisać jako
\begin{equation}
	s = \frac{\max(|x|)}{2^{b-1}-1}.
\end{equation}

Wariant asymetryczny wykorzystuje dodatkowo punkt zerowy $z_0$, który umożliwia przesunięcie zakresu kwantyzacji względem zera. Jest to przydatne w przypadku danych, których rozkład nie jest symetryczny. Wtedy dekwantyzacja przyjmuje postać
\begin{equation}
	\hat{x} = s(q - z_0).
\end{equation}

W modelach sieci neuronowych częściej stosuje się jednak kwantyzację symetryczną, ponieważ wagi po treningu zwykle mają rozkład zbliżony do symetrycznego wokół zera. Upraszcza to implementację \cite{frantar2023gptq}, ale potrafi wpłynąć negatywnie na rozmiar pliku i/lub zakres kwantyzacji ze względu na przechowywanie liczb ujemnych.

\subsection{Kwantyzacja per-tensor, per-channel i blokowa}

Sposób wyznaczania skali ma duże znaczenie dla jakości końcowego modelu. W kwantyzacji \textit{per-tensor} dla całego tensora stosuje się jeden wspólny współczynnik skali. Jest to rozwiązanie najprostsze i najtańsze obliczeniowo, ale może być zbyt mało precyzyjne dla dużych warstw o zróżnicowanych wartościach wag.

Dokładniejszym podejściem jest kwantyzacja \textit{per-channel}, w której osobną skalę wyznacza się dla każdego kanału, wiersza lub kolumny tensora, zależnie od implementacji. Taki wariant lepiej dopasowuje się do lokalnych różnic w rozkładzie wag, ale wymaga przechowywania większej liczby współczynników pomocniczych.

W praktycznych implementacjach bardzo często stosuje się także kwantyzację blokową, zwaną również grupową. W tym podejściu tensor dzieli się na mniejsze fragmenty o ustalonym rozmiarze, a dla każdego bloku wyznacza się osobną skalę. Pozwala to osiągnąć kompromis pomiędzy jakością rekonstrukcji a rozmiarem metadanych. Im mniejszy blok, tym dokładniejsze odwzorowanie wartości, ale tym większy narzut pamięciowy. Zależność ta będzie widoczna również w dalszych rozdziałach pracy, gdzie różne rozmiary bloków wpływają na końcowy rozmiar pliku oraz na jakość modelu po dekwantyzacji.

\subsection{Kwantyzacja po treningu i kwantyzacja podczas treningu}

Ze względu na moment wprowadzenia ograniczeń precyzji rozróżnia się dwa podstawowe podejścia. Kwantyzacja po treningu, czyli PTQ (\textit{Post-Training Quantization}), polega na skompresowaniu już wytrenowanego modelu bez ponownego uczenia wszystkich parametrów. Jest to rozwiązanie wygodne, ponieważ nie wymaga kosztownego procesu dostrajania modelu od początku, a przy odpowiednim doborze metody może zachować bardzo dobrą jakość \cite{frantar2023gptq}.

Drugi wariant to kwantyzacja podczas treningu, czyli QAT (\textit{Quantization-Aware Training}). W tym przypadku model jest uczony z uwzględnieniem ograniczeń wynikających z późniejszej niskiej precyzji zapisu. Pozwala to lepiej przygotować sieć na obecność błędów kwantyzacji, ale wymaga znacznie większych nakładów obliczeniowych i bardziej złożonego procesu trenowania.

W niniejszej pracy główny nacisk położony jest na podejście PTQ, ponieważ pozwala ono badać wpływ kompresji bez konieczności ponownego trenowania modeli. Jest to także podejście bliższe praktycznym narzędziom do kompresji modeli, stosowanym po zakończeniu ich uczenia.

\subsection{Błąd kwantyzacji}

Każde uproszczenie reprezentacji numerycznej powoduje błąd kwantyzacji. Można go rozumieć jako różnicę pomiędzy oryginalną wartością $x$ a jej odpowiednikiem po dekwantyzacji $\hat{x}$. Błąd ten można zapisać jako
\begin{equation}
	e = x - \hat{x}.
\end{equation}

W idealnym przypadku błędy te są niewielkie i rozkładają się w sposób, który nie wpływa znacząco na końcowy wynik modelu. W praktyce sytuacja zależy jednak od rozkładu wag, rodzaju warstwy oraz od tego, czy w danych występują wartości odstające. Właśnie dlatego niektóre modele zachowują bardzo dobrą jakość nawet po silnej kompresji, podczas gdy inne ulegają wyraźnej degradacji już przy umiarkowanym skróceniu precyzji.

Z punktu widzenia dalszych rozdziałów istotne jest więc przede wszystkim to, że skuteczna kwantyzacja nie polega wyłącznie na „ucięciu bitów”, lecz na takim dobraniu skali, sposobu grupowania i reprezentacji pomocniczej, aby zminimalizować wpływ błędu rekonstrukcji na zachowanie modelu.

\subsection{Notacja używana dalej w pracy}

W tym i dalszych rozdziałach pojawia się kilka grup symboli, które warto uporządkować już na tym etapie. Najważniejsze z nich związane są z kwantyzacją wag i z zapisem warstw modelu.

Oryginalną wartość liczbową oznaczam jako $x$, a jej wersję po kwantyzacji jako $q$. Wartość odtworzona w~procesie dekwantyzacji ma postać $\hat{x}$. Błąd kwantyzacji oznaczam symbolem $e$. Funkcję kwantyzacji zapisuję jako $Q$, a funkcję dekwantyzacji jako $D$.

W przypadku wag modeli przyjmuję zapis $w$ dla pojedynczej wagi oraz $W$ dla całej macierzy wag. Skwantyzowany odpowiednik macierzy wag oznaczam jako $\hat{W}$, natomiast macierz wejść jako $X$. W rozdziale o GPTQ pojawia się również $H_F^{-1}$, czyli odwrotna macierz Hesjana, wykorzystywana przy analizie wpływu błędu na wyjście warstwy. Gdy w tekście mowa o różnicy pomiędzy wartością oryginalną i przybliżoną w bardziej ogólnym sensie, używany jest symbol $\Delta$.

Przy opisie skali i przesunięcia stosuję symbole $s$, $z$ oraz $z_0$. W zależności od kontekstu $s$ oznacza współczynnik skali używany do odwzorowania liczb z oryginalnej domeny na domenę skwantyzowaną i z powrotem. Symbol $z$~oznacza przesunięcie lub punkt zerowy w zapisie ogólnym, natomiast $z_0$ jest punktem zerowym w zapisie bardziej precyzyjnym, szczególnie wtedy, gdy trzeba wyróżnić wariant asymetryczny.

W częściach dotyczących blokowej kompresji wag pojawiają się także symbole $b$ i $B$. Litera $b$ oznacza liczbę bitów przeznaczonych na zapis jednej wartości, a $B$ rozmiar bloku, czyli liczbę wag grupowanych wspólnie przy wyznaczaniu skali i innych parametrów pomocniczych. Jeżeli tekst odnosi się do granic reprezentacji dla danego formatu, pojawia się również zapis zakresu dopuszczalnych wartości, np. dla INT4 lub INT8.

W opisie warstw i modeli często występują też zapisy związane z architekturą sieci. Dense, Conv1D, Conv2D, GRU oraz LayerNormalization oznaczają konkretne typy warstw, a GeLU, Sigmoid i Softmax oznaczają funkcje aktywacji używane w ich wyjściach. Gdy w tekście mowa o jakości modelu, stosowane są nazwy metryk: dokładność, precyzja, czułość, miara F1, MAE, MSE i RMSE.

W częściach dotyczących implementacji i testów pojawiają się ponadto skróty obliczeniowe i organizacyjne. CPU oznacza procesor centralny, GPU procesor graficzny, CUDA platformę obliczeń równoległych, a kernel CUDA pojedynczą funkcję wykonywaną równolegle na GPU. Pojęcie \textit{batch size} oznacza rozmiar paczki danych przetwarzanych w jednym kroku, a \textit{block size} rozmiar bloku wag wykorzystywanego w kwantyzacji.

\chapter{Popularne rozwiązania}
\chapter{Popularne metody kwantyzacji modeli}

Rozwój dużych modeli językowych (LLM) oraz modeli wizji komputerowej (VLM) doprowadził do sytuacji, w której wdrożenia produkcyjne oraz badania nad architekturami transformerów stały się domeną wyłącznie korporacji dysponujących wydajnymi klastrami obliczeniowymi. Kwantyzacja, polegająca na redukcji precyzji reprezentacji liczb (najczęściej z 32-bitowego formatu zmiennoprzecinkowego do 4- lub 8-bitowych wartości stałoprzecinkowych bądź niestandardowych formatów zmiennoprzecinkowych), stała się kluczowym mechanizmem zwiększającym popularność sztucznej inteligencji. Pozwala ona na drastyczne zmniejszenie zapotrzebowania na przepustowość pamięci oraz pojemność pamięci, przyspieszając wnioskowanie i umożliwiając uruchamianie modeli na sprzęcie konsumenckim.

W niniejszym rozdziale przedstawione zostaną najbardziej rozpowszechnione rozwiązania dziedzinie kwantyzacji. Omówiono zarówno standardy  BitsAndBytes z biblioteki Hugging Face, algorytmy kompresji wag bazujące na analizie błędów (AWQ, GPTQ), nowatorskie formaty plików zoptymalizowane pod heterogeniczne architektury CPU i GPU (GGUF), jak i przełomowe metody mapowania rozkładu wag oparte na transformacjach ortogonalnych (TurboQuant).

\section{BitsAndBytes oraz format NF4}
\label{sec:bnb}

Choć biblioteka BitsAndBytes (BnB) początkowo funkcjonowała jedynie jako biblioteka w ekosystemie PyTorch, dziś nazwa ta utożsamiana jest z de facto standardem w kwantyzacji, głównie dzięki artykułowi \textit{QLoRA} (Dettmers i in., 2023), który spopularyzował format \textbf{NF4} (NormalFloat4). Metoda ta pozwoliła na dostrojenie modeli rzędu 65 miliardów parametrów na pojedynczym akceleratorze graficznym, co wcześniej było niemożliwe.

Standardowe kwantyzacje (np. INT8) zakładają jednostajny rozkład wartości. Oznacza to, że odległość między poszczególnymi poziomami reprezentacji (krok kwantyzacji $\Delta$) jest stały. Badania wykazały jednak, że wytrenowane wagi modeli sieci neuronowych przyjmują rozkład zbliżony do normalnego (o zerowej średniej, skupiony wokół zera z długimi ogonami). W związku z tym przypisanie równych fragmentów pamięci bitowej dla rzadkich wartości na ogonach i gęsto skupionych wartości wokół zera prowadzi do niepotrzebnej straty precyzji (tzw. błąd kwantyzacji).

BitsAndBytes rozwiązuje to za pomocą \textbf{kwantyzacji kwantylowej}. Algorytm wyznacza wartości odczytane z funkcji odwrotnej do skumulowanej funkcji rozkładu normalnego (z ang. Inverse Cumulative Distribution Function - ICDF). W ten sposób uzyskuje się optymalne, \textbf{niestandardowe poziomy wartości}, które gwarantują, że każdy z tzw. "przedziałów" (z ang. bins) posiada dokładnie takie samo prawdopodobieństwo wystąpienia wagi z danego przedziału.

\subsection{Warianty i podtypy w architekturze BitsAndBytes}

Ekosystem BitsAndBytes rozwinął się w stronę wsparcia wielu architektur i potrzeb treningowych, wprowadzając szereg podtypów kwantyzacji:

\begin{itemize}
    \item \textbf{LLM.int8()}: Pionierska metoda umożliwiająca wnioskowanie (inferencję) w formacie 8-bit bez zauważalnego spadku wydajności. Problemem wczesnych modeli była anizotropia aktywacji – podczas wnioskowania pojawiały się tzw. wartości odstające w stanach ukrytych, które mogły mieć amplitudę nawet 100 razy większą niż reszta wymiarów. Standardowa kwantyzacja INT8 niszczyła model. LLM.int8() wprowadziło hybrydową dekompozycję: macierze mnożone są w formacie INT8 z funkcją typu mieszanej precyzji (z ang. "mixed-precision"). Kanały zawierające wartości odstające (zazwyczaj ok. 0.1\% wszystkich wymiarów) izolowane są do mnożenia w pełnej precyzji FP16.
    
    \item \textbf{NF4 (NormalFloat4)}: Główny format używany do kompresji pamięciowej podczas treningu z wykorzystaniem techniki QLoRA. Należy tu podkreślić, że dekwantyzacja z NF4 do formatu BF16/FP16 na potrzeby obliczeń zachodzi w trakcie działania. Oszczędność wynika z faktu, że wagi modelu spoczywające w pamięci VRAM zajmują ułamek oryginalnej wielkości.
    
    \item \textbf{FP4 (Float4)}: Alternatywny format dla NF4. Zamiast niestandardowych przedziałów mieszczących się w rozkładzie normalnym, używa klasycznej reprezentacji zmiennoprzecinkowej z bitem znaku, wykładnikiem i mantysą. Ze względu na bardzo małą liczbę bitów wykładnika, jest on niezwykle podatny na tzw. niedomiar/przepełnienie i jest zazwyczaj gorzej oceniany w testach niż NF4.
    
    \item \textbf{Podwójna Kwantyzacja (Double Quantization - DQ)}: Nawet przy kwantyzacji do 4 bitów potrzebne są stałe normalizujące (skalowanie i bias) zazwyczaj przechowywane w formacie FP32 dla każdej grupy wag. BnB wprowadziło podwójną kwantyzację, która kwantyzuje same stałe (tzw. metadane) z FP32 do FP8, co pozwala zaoszczędzić dodatkowe 0.37 bita na każdą wagę oryginalnego modelu. Co ciekawe, w BnB DQ parametry pomocnicze rozmiaru grupy wynoszą 256, co zachowuje strukturalną integralność macierzy Hesjana procesu uczenia.
\end{itemize}

\section{GPTQ (Accurate Post-Training Quantization)}
\label{sec:gptq}

Algorytm GPTQ, zaproponowany przez Frantara pod koniec 2022 roku, wyznaczył zupełnie nowe standardy dla kwantyzacji po treningu (PTQ - Post-Training Quantization). Jego główną zaletą jest to, że potrafi on zredukować precyzję modeli posiadających setki miliardów parametrów (np. modele takie jak OPT-175B czy BLOOM-176B) w zaledwie kilka godzin obliczeń GPU, przy zachowaniu niemal identycznej dokładności względem modelu oryginalnego.

GPTQ bazuje na koncepcjach algorytmu OBD (Optimal Brain Damage) z początku lat 90., który służył do "przycinania" sieci neuronowych. W skrócie, GPTQ rozwiązuje następujący problem: mając warstwę liniową o wagach $W$, chcemy znaleźć skwantyzowane wagi $\hat{W}$, tak aby błąd rekonstrukcji wyjścia warstwy (błąd aktywacji) był jak najmniejszy dla danego wejścia $X$. Minimalizowana jest funkcja:
\begin{equation}
    \arg\min_{\hat{W}} || WX - \hat{W}X ||^2
\end{equation}

Ideą algorytmu jest kompresja wag wierszami (lub kolumnami w zależności od implementacji) oraz obserwowanie błędu powstającego na wyjściu, by ten sam błąd natychmiast kompensować na wagach sąsiednich, jeszcze nie skwantyzowanych. Obliczanie optymalnych modyfikacji wymaga zróżniczkowania funkcji błędu po wadze względem pozostałych wag z wykorzystaniem Macierzy Hesjana $H_F^{-1}$.

Aby algorytm zadziałał prawidłowo, macierz Hesjana jest statystycznie aproksymowana na małym zbiorze kalibracyjnym – wymagane jest zaledwie kilkaset próbek danych wejściowych. Aby zapobiec załamaniu numerycznemu (inwersja macierzy dla tysięcy wymiarów jest niestabilna), twórcy wdrożyli szybką dekompozycję Cholesky'ego i zjawisko "opóźnionych aktualizacji partii". GPTQ wylicza korekcję w paczkach po 128 kolumn jednocześnie.

\subsection{Warianty i podtypy w ekosystemie GPTQ}

Podczas gdy sam algorytm bazowy GPTQ jest tylko matematycznym modelem PTQ, społeczność open source (zwłaszcza projekt AutoGPTQ) wypracowała cały system podtypów i wariantów dopasowanych do sprzętu.

\begin{itemize}
    \item \textbf{Różne szerokości bitowe}: Najpopularniejszym formatem zaimplementowanym w GPTQ jest kompresja do \textbf{4-bitowych liczb całkowitych (INT4)}. Oferuje on najlepszy stosunek kompresji do spadku "perplexity". Warianty takie jak INT3 i INT8 również są dostępne, jednak INT3 powoduje już zauważalną degradację właściwości logicznych, a INT8 daje zbyt małą kompresję, by uzasadnić narzut dekwantyzacji.
    
    \item \textbf{Rozmiar Grupy}: GPTQ stosuje kwantyzację grupową. Aby równanie kompensacji Hesjana nie okazało się zbyt ogólne (co ma miejsce w kanałowej), a jednocześnie nie pochłaniało nadmiernie pamięci na metadane (jak w przypadku elementowej), grupuje się wagi w bloki. Najpopularniejsze wartości parametru rozmiar grupy to:
    \begin{itemize}
        \item \textbf{-1 (Kanałowa/Kolumnowa)}: Jeden zestaw parametrów (skala, punkt zerowy) dla całego wymiaru wejściowego. Oszczędność pamięci jest największa, ale model jest podatny na anomalię dzielenia przez zero.
        \item \textbf{128}: Standardowa kompresja zbalansowana. Każde 128 wag posiada własny współczynnik skali $s$ i przesunięcie $z$.
        \item \textbf{32}: Bardzo wysoka precyzja dekwantyzacji. Metadane dla grup 32-bitowych potrafią zająć tyle samo miejsca co skwantyzowane wagi, co w rezultacie daje bitrate na poziomie 5.5-6 bitów zamiast założonych 4.
    \end{itemize}
    
    \item \textbf{desc-act (Kolejność aktywacji / statystyki rozkładu GPTQ)}: Znaczący podtyp algorytmu. Badacze zauważyli, że jeśli algorytm nie kwantyzuje macierzy losowo lub sekwencyjnie, ale w kolejności malejącej wariancji aktywacji (malejącej istotności macierzy Hesjana), to błędy zaokrągleń mogą zostać drastycznie zredukowane. Ten wariant określa się skrótowo jako \textbf{desc-act} (zazwyczaj w nazwie pliku modelu). Zwiększa on jednak czas kompresji dwu-trzykrotnie.
\end{itemize}

\section{AWQ (Activation-Aware Weight Quantization)}
\label{sec:awq}

Publicacja AWQ (Activation-Aware Weight Quantization) z 2023 roku wniosła do dziedziny nową perspektywę: perspektywę braku symetrii ważności poszczególnych kanałów modelu. Zauważono, że do tej pory metody (w tym GPTQ) analizowały błędy z perspektywy globalnej w przestrzeni wag modelu. AWQ dowodzi, że w architekturze modelu istnieje zaledwie około \textbf{0,1\%-1\% wag} które mają kluczowe znaczenie w procesie wnioskowania i modyfikowanie ich wywołuje katastrofalny spadek inteligencji modelu.

Te wagi nadrzędne są ściśle skorelowane ze zjawiskiem odpowiadającym im aktywacji o dużej amplitudzie. W przeciwieństwie do aktywacji odstających w BnB, w badaniach AWQ wykazano, że są to stałe cechy konkretnych kanałów, charakterystyczne dla typowych dystrybucji tekstowych; kanał, czy kolumna macierzy $W$ - otrzymuje zawsze silne pobudzenia w wektorze wejściowym $X$.

Podstawową różnicą AWQ względem GPTQ jest podejście do kalibracji. O ile ten pierwszy szuka optymalnego składnika błędu Hesjanowego, AWQ nie komplikuje sobie życia zaawansowaną algebrą liniową na poziomie szkolenia. Zamiast tego AWQ zauważa, że jeśli chcemy zapisać istotną wagę $w$ – to nawet duży błąd kwantyzacji $\Delta w$ przy niskiej intensywności aktywacji $x$ nie wywoła wielkich szkód w iloczynie wyjściowym. Jednakże – ten sam błąd wielkości $\Delta w$ przy bardzo dużej aktywacji $x$ spowoduje zniekształcenie wyjścia $x \cdot \Delta w$ całkowicie degradujące wynik. AWQ rozwiązuje to zjawisko stosując proste funkcjonalne \textbf{skalowanie (per-channel scaling)}. Tworzy wektor $s$ oraz szuka argumentu minimalizującego błąd replikacji $\mathrm{round}(w/s)$.

Dzięki takiemu przemyśleniu problemu, metoda AWQ potrafi drastycznie kompresować model analogicznie do GPTQ, jednak zachowując spójność dla tych 1% newralgicznych strumieni danych.

\subsection{Podtypy, warianty sprzętowe i architektoniczne}

Podobnie jak i GPTQ, AWQ rozwijane było jako projekt o otwartym kodzie źródłowym. Wyróżniamy w nim szereg wariantów wykonawczych i podtypów:

\begin{itemize}
    \item \textbf{Warianty bitowe W$X$A$Y$}: Architektury kwantyzowane AWQ są najczęściej opisywane notacją WxAy. Odnosi się to do precyzji wag oraz aktywacji. Bazowym i najpopularniejszym podtypem jest \textbf{W4A16} (4-bitowe wagi, 16-bitowe aktywacje). Autorzy zaproponowali również \textbf{W3A16} oraz eksperymentalne architektury \textbf{W8A8}. Ten wariant dotyczy problematycznej kwantyzacji samych wektorów wejściowych co jest obarczone bardzo dużym błędem maksymalnej wartości bezwzględnej.
    
    \item \textbf{Warianty kerneli (GEMM vs GEMV)}: Bardzo ważnym podziałem AWQ jest klasyfikacja ze względu na maszyny tensorowe na docelowym sprzęcie GPU. Są to tzw. Kernels w bibliotece AWQ:
    \begin{itemize}
        \item \textbf{GEMM-based (GEMM Kernel)}: Mnożenie macierzy zoptymalizowane pod pętle w architekturach Ampere (NVIDIA RTX 3000) i Hopper. AWQ zamiast dekwantyzować wagi na bieżąco do FP16, przekształca wagi do postaci rozpakowanej tuż przed wejściem do Tensor Core, opierając się o wzorce przekształceń GEMV. Format zalecany przy maksymalizacji liczby klatek na sekundę dla dużych rozmiarów partii (uruchamianie autorskich serwerów vLLM z włączonym stronnicowaniem).
        \item \textbf{GEMV-based (GEMV Kernel)}: General Matrix-Vector. Klasyczne dekwantyzacje przeprowadzane wyłącznie przy użyciu rejestrów CUDA. Stworzone specjalnie dla publikacji TinyChat i kart graficznych o bardzo małym standardzie obliczeniowym, takich jak mobilne układy Nvidia Jetson. Profile te są nastawione na wnioskowanie z rozmiarem partii, batch\_size=1. W ekosystemie jest to tzw. przypadek graniczny alokacji.
    \end{itemize}
    \item \textbf{Zoptymalizowane rozmiary grupy}: Podtypy AWQ zazwyczaj przyjmują wielkość bloku równą 128 (podobnie jak było to zaprezentowane w AutoGPTQ), co gwarantuje optymalny design na urządzeniach komórkowych. Wiele implementacji sprzętowych rezerwuje rozmiar grupy na 32 dla profilu wysokich architektur pro-użytkowych (np. procesory NPU).
\end{itemize}

\section{GGUF (GPT-Generated Unified Format)}
\label{sec:gguf}

Format GGUF (GPT-Generated Unified Format) to podstawowy fundament współczesnej inżynierii modeli na urządzenia brzegowe. Zaproponowany przez Georgiego Gerganowa (twórcę oprogramowania \texttt{llama.cpp}), wyparł on w szybkim czasie starszy format GGML. Motywacją do powstania GGUF był brak kompatybilności wstecznej i trudności w rozszerzaniu starych formatów o nowe architektury sieci, a także brak natywnego wsparcia dla zaawansowanych metod kwantyzacji.

GGUF definiuje strukturę przechowywania metadanych w eleganckiej strukturze typu klucz-wartość (KV) na początku pliku przechowującej parametry modelu, a tuż za nimi mapę wielkich tensorów reprezentującą wytrenowane, skwantyzowane parametry. Jest to format zoptymalizowany na poziomie C/C++ z natychmiastowym wsparciem operacji mapowania pamięci bez kopiowania (mmap), co ułatwia ładowanie modeli do pamięci RAM bez konwersji oraz wykorzystuje instrukcje SIMD (AVX2, AVX-512 na x86, NEON na architekturze ARM).

W przeciwieństwie do poprzednich metod (operujących na pojedynczych macierzach, jak AWQ), GGUF z założenia operuje na kompromisach sprzętowych dla poszczególnych warstw - to znaczy pozwala na ładowanie konkretnych warstw typu attention w różnej rozdzielczości w zależności od zapotrzebowania procesora.

\subsection{Architektura i mapa wartości GGUF}

Format GGUF jest niezwykle złożony, a wyselekcjonowane typy danych są indeksowane w sposób ciągły w celu łatwej interpretacji przez parser C++. Macierz pełnej mapy tensorów w formacie GGUF przedstawia się następująco:

\begin{table}[h]
\centering
\caption{Pełna mapa typów kwantyzacji w formacie GGUF}
\label{tab:gguf_map_full}
\begin{tabular}{|c|l|c|l|}
\hline
\textbf{Indeks} & \textbf{Format} & \textbf{Indeks} & \textbf{Format} \\ \hline
0 & F32 & 19 & IQ2\_XSS \\ \hline
1 & F16 & 20 & IQ2\_XS \\ \hline
2 & Q4\_0 & 21 & Q2\_K\_S \\ \hline
3 & Q4\_1 & 22 & IQ3\_XS \\ \hline
4 & Q4\_1\_SOME\_F16 & 23 & IQ3\_XXS \\ \hline
5 & Q4\_2 & 24 & IQ1\_S \\ \hline
6 & Q4\_3 & 25 & IQ4\_NL \\ \hline
7 & Q8\_0 & 26 & IQ3\_S \\ \hline
8 & Q5\_0 & 27 & IQ3\_M \\ \hline
9 & Q5\_1 & 28 & IQ2\_S \\ \hline
10 & Q2\_K & 29 & IQ2\_M \\ \hline
11 & Q3\_K\_S & 30 & IQ4\_XS \\ \hline
12 & Q3\_K\_M & 31 & IQ1\_M \\ \hline
13 & Q3\_K\_L & 32 & BF16 \\ \hline
14 & Q4\_K\_S & 33 & Q4\_0\_4\_4 \\ \hline
15 & Q4\_K\_M & 34 & Q4\_0\_4\_8 \\ \hline
16 & Q5\_K\_S & 35 & Q4\_0\_8\_8 \\ \hline
17 & Q5\_K\_M & 36 & TQ1\_0 \\ \hline
18 & Q6\_K & 37 & TQ2\_0 \\ \hline
 &  & 38 & MXFP4\_MOE \\ \hline
\end{tabular}
\end{table}

W tabeli tej można wyróżnić kilka rodzin rozwiązań, które często sprawiają trudności interpretacyjne. Powszechnie oznaczenia cyfrowe na pierwszej pozycji (\textit{Qx}) odnoszą się do docelowej pojemności kontenera (bity na wagę - bpw). Oznaczenia \textbf{"\_0" (np. Q4\_0, Q5\_0)} odnoszą się do modeli \textbf{z bezpośrednim skalowaniem}. Oznaczenia \textbf{"\_1" (np. Q4\_1, Q5\_1)} posiadają dodatkowo odchylenie minimum (tzw. przesunięcie (offset) dla każdego bloku, co daje większą rozdzielczość dynamiczną). Z kolei nowością sprzętową są podtypy indeksowe 33-35 (np. \textbf{Q4\_0\_4\_4}), które definiują modele określane jako "sub-block matrix" i dedykowane są one układom NPU w architekturze ARM gdzie instrukcje SIMD optymalizowane są na poziomie sprzętu w alokacjach pamięci L1/L2 4x4, 4x8 i 8x8 bajtów.

\begin{table}[h]
    \centering
    \caption{Wpływ kwantyzacji GGUF na rozmiar modelu na przykładzie Qwen3.6-27B}
    \label{tab:gguf_size_impact}
    \begin{tabular}{|c|c|}
        \hline
        \textbf{Format} & \textbf{Rozmiar Modelu (GB)} \\ \hline
        F16 & 55.6 \\ \hline
        IQ4\_XS & 15.4 \\ \hline
        IQ4\_NL & 16.1 \\ \hline
        Q3\_K\_S & 12.4 \\ \hline
        Q3\_K\_M & 13.6 \\ \hline
        Q4\_0 & 15.8 \\ \hline
        Q4\_1 & 17.3 \\ \hline
        Q4\_K\_S & 15.9 \\ \hline
        Q4\_K\_M & 16.8 \\ \hline
        Q5\_K\_S & 19.2 \\ \hline
        Q5\_K\_M & 19.5 \\ \hline
        Q6\_K & 22.5 \\ \hline
        Q8\_0 & 28.6 \\ \hline
    \end{tabular}
\end{table}


\subsection{Kwantyzacja typu K-Quants}
\textit{K-Quants} to zaimplementowana w GGUF metoda kwantyzacji podrzędnej, niejednorodnej na poziomie całego modelu. Bazuje na aproksymacji ważności tensorów w podziale na trzy poziomy wrażliwości: mały (Ma), średni (Mb) i duży (Mc). Rozwiązanie \textbf{K\_S (Small)} stosuje "najmocniejszą", najbardziej drastyczną formę redukcji. Generowane są tu profile Q2\_K, Q3\_K\_S.
Stosując podejście skwantyzowane do macierzy attention, okazało się, że tzw. mapa gradientowa jest niejednorodna, przez co powstały kolejne podrodziny:
\begin{itemize}
    \item \textbf{K\_M (Medium)}: Uznany zazwyczaj za kompromis rozwiązań progresywnych najmocniej kompresujących, a standardowym profilem. W przypadku podtypu Q4\_K\_M warstwy mniej istotnych tensorów oraz wyjść modelu są kompresowane w 4 bitach, ale newralgiczne bloki w modelach są kompresowane w wyższej rozdzielczości (6 bitów).
    \item \textbf{K\_L (Large)}: Profile z minimalnym uszczerbkiem na jakości. Utrzymują większość podmatryc na poziomie 6 lub 8 bitów, a kompresji do 4 bitów poddawane są jedynie ogony tensorów.
\end{itemize}

\subsection{Importance-Matrix oraz rodzina IQ}
Prawdziwym przełomem inżynierii wstecznej w metodzie GGUF jest \textbf{Importance Matrix} i podtypy \textbf{IQ} (Importance Quantization). W tej metodzie nie określamy istotności prosto z "wariancji aktywacji" lecz statystycznie z kalibracyjnego zbioru treningowego wyliczamy entropię punktową. Otrzymujemy n+k-wymiarowy tensor "ważności" dla wag, bazując na uogólnionym błędzie dowolnym.
Formaty takie jak \textbf{IQ1\_S, IQ1\_M} (indeksy 24, 31) umożliwiają agresywną kompresję do 1 bita na wagę (przez co wielkość pliku maleje ponad 16-krotnie!). Metody te są często znacznie lepsze od współczesnych podejść z rodziny klasycznego Q2 i Q3, właśnie ze względu na zgodność informacyjną z macierzą ważności.
Ostatnie aktualizacje GGUF z 2024 r. wprowadzają wsparcie dla podtypów architektur mixture-of-experts (MoE), takich jak \textbf{MXFP4\_MOE} (indeks 38) optymalizujących alokację routingów na poziomie bramek ekspertów.

\section{TurboQuant}
\label{sec:turboquant}

Na tle powyższych, wysoce rozwiniętych i sprofilowanych rozwiązań, TurboQuant stanowi nowatorskie rozwiązanie, które zdołało całkowicie namieszać w świecie optymizacji. Jest to metoda wywodząca się z dziedziny badań nad optymalizacją modeli i charakteryzuje się hybrydowym podejściem.

Problematyka operacji na liczbach całkowitych zakłada, że błędy skwantyzowanych prążków zwiększają się liniowo (zwiększa stratę sVRG) dla macierzy Hesjana wpisującego się w założenia rozszerzonych sieci rezydualnych, co przy modelach o skali 7B-70B powoduje tzw. cyfrową degradację wyjść. TurboQuant, zamiast korzystać ze skomplikowanej kalibracji z użyciem Hessiana jak GPTQ i AWQ, bezpośrednio zanurza wielowymiarową przestrzeń tensorów rzadkich do przestrzeni "gładkich" za sprawą transformacji Hadamarda i losowych macierzy ortogonalnych.

Główną ideą TurboQuant jest zaobserwowanie, że wstępnie skwantyzowana macierz wejścia według normy L2 w architekturze LLM ma strukturę quasi-izotropową. Zamanifestowało się to w filozofii losowych rotacji. Klastrowanie macierzy ortogonalnych i wykonanie tzw. szybkiej transformacji Walsha-Hadamarda (Fast Walsh-Hadamard Transform - FWHT) dla wejść przed włączeniem przybliżenia całkowitego, wytraca wagom ich wartości odstające w sposób prawie idealny. Z matematycznego punktu widzenia, TurboQuant nadaje macierzom "izotropię energetyczną" – dystrybucja jest spłaszczana. To usuwa anizotropię która prowadziła do dotychczasowych błędów. Fizycznie, usunięcie dużych oscylacji powoduje, że proste zaokrąglenie do INT4 czy nawet INT3 nie powoduje zapadania się klastrów odpowiadających za rozumowanie logiczne (Chain-of-Thought).

\subsection{Podtypy TurboQuant}

???

\begin{itemize}
    \item \textbf{TQ1\_0 oraz TQ2\_0}: Skwantyzowane z rzędu wektory istotnych cech zaimplementowanych do struktury GGUF (patrz Tabela \ref{tab:gguf_map_full}). W przeciwieństwie do klasycznych kwantyzacji 1-bitowych (jak binarne BNN), TQ2\_0 dzięki losowej rotacji gwarantuje, że w przestrzeni 2 bitów błąd kwantyzacji jest wysoce nieproporcjonalnie pomniejszony. Wdrożenia pokazały spadek jedynie o 1-2\% w testach MMLU i HumanEval w przeciwieństwie do katastrofalnych 20\% w starszych metodach 2-bitowych. Metoda TQ1\_0 jest absolutnym ekstremum kompresji, zaledwie 1 bit na wagę (wymaga dedykowanych kerneli dekompresujących).
    
    \item \textbf{TurboQuant-A (Approximate)}: Oparty na ucięciu (pruning) mało istotnych wektorów bazowych i transformacji Hadamarda w blokach 256x256. Umożliwia to wnioskowanie w czasie rzędu O(N log N).
    
    \item \textbf{Integracja z MoE}: Z racji wielkich osiągnięć TQ w utrzymaniu gładkiego rozkładu, metoda ta posłużyła jako fundament dla współczesnych MoE (Mixture of Experts) gdzie tradycyjny błąd kwantyzacji mógł zrujnować systemy routingu ekspertów. Podtyp \textbf{TQ-MoE} tworzy klastrową rotację dla ekspertów w siatce nawiązując bezpośrednio do formatów takich jak MXFP4\_MOE (indeks 38 w mapie GGUF). Pokazano tam, że TQ-MoE jest w stanie skompresować 8x7B model ekspertowy do reprezentacji 2-bitowej z lepszymi wynikami niż AWQ w 4-bit dla oryginalnych rozmiarów.
\end{itemize}

Ze względu na to, jak nowe rozwiązanie znacząco zmieniło kompresję (minimalizując wymóg kalibracji zbioru uczącego do niewielkiej liczby próbek, a często kalibrując w trakcie działania bez zbioru danych - zero-shot), TurboQuant stanowi dziś bardzo ważny punkt odniesienia w świecie vLLM i CUDA. Jego hybrydowe implementacje wymusiły na twórcach optymalizacje zunifikowane dla AVX-512.
\chapter{Własna implementacja}
\section{Opis implementacji}
Badając metody kwantyzacji modeli sztucznej inteligencji natrafiłem na wzór prostej kompresji wag. Zdecydowałem się skorzystać z niego w ramach tej pracy licencjackiej i stworzyć cały plik w stylu biblioteki w języku Python do szybkiej kwantyzacji i dekwantyzacji wag modeli tensorflow.
\subsection{Kwantyzacja}
Skrótowa, punktowa prezentacja procesu kwantyzacji::
\begin{itemize}
    \item Wczytywane są warstwy modelu i następuje iteracja po nich,
    \item Dla każdej warstwy odczytywane są jej wagi,
    \item Wagi dzielone są na bloki o określonym rozmiarze (np. 32 wagi),
    \item Dla każdego bloku obliczana jest jej skala według wzoru,
    \item Każda waga bloku jest dzielona przez skalę i zaokrąglana do najbliższej liczby całkowitej oraz "ściskana" do ustalonego zakresu wartości,
    \item Takie wagi są przesuwane do pozytywnego zakresu wartości i zapisywane do pliku w postaci binarnej.
\end{itemize}

Kwantyzacja zaczyna się od pobrania informacji o warstwach modelu i iteracji po ich wagach. Przy każdej warstwie sprawdzana jest liczba jej wymiarów, obsługiwane są warstwy o 1-, 2-, 3- i 4-wymiarowe. Każdy z tych wariantów ma swoją własną ścieżkę w kodzie, ponieważ większa liczba wymiarów wymaga zagnieżdżenia większej liczby pętli.
Wewnątrz tych pętli wagi są dzielone na bloki o określonym rozmiarze (np. w swoich testach korzystałem z takich wartości dla partii jak 8, 16, 32, 64 i 128). Do kompresji dla każdego bloku wyliczany jest jego współczynnik skalowania według wzoru:
$\mathit{half\_point} = \frac{2^{\mathit{bits}}}{2}$
$\mathit{scale} = \frac{\max\!\big(|\mathit{block}|\big)}{\mathit{half\_point} - 1}$
Gdzie $\mathit{block}$ to zbiór wag w danym bloku, a $\mathit{bits}$ to liczba bitów, na które chcemy skompresować wagi (kod obsługuje opcje 4 i 8 bit).
Mając skalę, każda waga bloku jest dzielona przez nią i zaokrąglana do najbliższej liczby całkowitej. Następnie jest "ściskana" do zakresu wartości dyktowanego przez liczbę bitów, czyli $[-\mathit{half\_point} - 1, \; \mathit{half\_point} - 1]$.
Na koniec, aby nie poświęcać jednego bitu na dodatkowe informacje, $\mathit{half\_point}$ jest dodawane do każdej wagi, przesuwając je do pozytywnego zakresu wartości.
Takie wagi są następnie zapisywane do pliku w postaci binarnej wraz ze skalą. Zapis następuje w formacie: [skala (float32), wagi (uint8)]. Uint8 jest stosowany nawet w przypadku opcji 4-bit, a wtedy do na jedną liczbę uint8 składają się dwie wagi 4-bitowe złożone z sobą ($\mathit{w_1}\text{ << 4 || }\mathit{w_2}$).
Do pliku binarnego nie trafiają żadne dodatkowe dane i nie są wstawiane żadne separatory między warstwami. Informacje o strukturze modelu, takie jak liczba warstw, ich rozmiary i rodzje, są zachowywane poprzez zapis obok pliku binarnego pliku json, korzystając z wbudowanych funkcji eksportu konfiguracji modelu tensorflow. Dodatkowe informacje, takie jak rozmiar bloku i liczba bitów, są zapisywane do osobnego pliku json.
Na dysku te trzy pliki nie występują oddzielnie, a są zapisywane do jednego pliku z niestandardowym rozszerzeniem \textit{.uniq}, którego nazwa pochodzi od nazwy wybranej dla biblioteki, \textit{Uniquant}. Plik ten jest w rzeczywistości plikiem typu archiwum. Dzięki temu wszystkie wymagane dane są przechowywane w jednym miejscu, co ułatwia przechowywanie oraz dekwantyzację.

Warto też wspomnieć, że w przypadku, gdy dana warstwa na ostatnim swoim wymiarze ma mniej wag niż określony rozmiar bloku, to jest zapisywana w pełnej postaci float32 bez kompresji. Zostało to dodane z dwóch powodów. Po pierwsze, ułatwia to dekompresję takich warstw. Po drugie, znacznie zmniejsza to wpływ na kompresji na jakość wynikowych wag po dekompresji.

\subsection{Dekwantyzacja}
Proces dekwantyzacji jest odwrotnością kwantyzacji, ale jest od niej o wiele złożony. Większa trudność wynika z faktu, że kwantyzacja tylko iteruje po kolejnych wagach, a granice warstw są nieprzekraczalne, podczas gdy w dekwantyzacji odczytywane wagi są zapisane w ciągłym wektorze i należy odpowiednio wyliczyć ich granice, odróżnić skale od wag i przypisać je w odpowiednim formacie do odpowiednich warstw modelu.
Skrótowa, punktowa prezentacja procesu dekwantyzacji:
\begin{itemize}
    \item Odczytywane są informacje kwantyzacji z jednego pliku json, a informacje o strukturze modelu z drugiego,
    \item Odczytywany jest plik binarny z wagami i skalami,
    \item Następuje iteracja po warstwach modelu,
    \item Wydzielany jest zakres wag warstwy,
    \item Wagi są mnożone przez odpowiednie skale i cofane jest ich przesunięcie,
    \item Takie zdekwantyzowane wagi są przypisywane do odpowiedniej warstwy modelu.
\end{itemize}

Dekwantyzacja zaczyna się od odczytania informacji o kwantyzacji (rozmiar bloku, liczba bitów) z jednego pliku json z paczki \textit{.uniq} oraz informacji o strukturze modelu z drugiego pliku json. Te drugie wczytywane są funkcjami tensorflow, które odtwarzają z nich model z pustymi wagami. Następnie wczytywany jest plik binarny z wagami i skalami do zmiennej, która będzie przechowywać te wagi przez cały proces dekwantyzacji.
Po tym rozpoczyna się pętla po warstwach modelu w postaci iteracji po obiektach struktury modelu. Po stronie dekwantyzacji ścieżki kodu nie są podzielone na podstawie liczby wymiarów warstw, a na podstawie nazw typów warstw. Obsługiwane są warstwy Dense, Conv1D, Conv2D, GRU i LayerNormalization. Na początku kodu dla każdej warstwy jest wyliczany zakres jej danych w wektorze wag na podstawie rozmiarów warstwy oraz wskaźnika aktualnej pozycji w wektorze wag, który jest aktualizowany po każdej warstwie.
$\mathit{w}[\mathit{ptr} : \mathit{ptr} + \mathit{layer\_size}]$
Gdzie $\mathit{w}$ to wektor wag, $\mathit{ptr}$ to wskaźnik aktualnej pozycji, a $\mathit{layer\_size}$ to liczba bitów danych warstwy (współczynniki skali i wagi). Ta ostatnia zmienna rozmiaru jest wyliczana poprzez pomnożenie razy 4 (lub 8, w zależności od parametru bitów kwantyzacji) liczby wszystkich wag warstwy oraz dodanie do tego pomnożonej przez 32 (liczba bitów reprezentacji skali) liczby bloków warstwy.
Ponieważ warstwy tensorflow składają się w rzeczywistości z 2 lub więcej macierzy wag, należy brać to pod uwage w trakcie wyliczania zakresu danych warstwy oraz podczas faktycznej dekwantyzacji wag. Liczba macierzy określa liczbę głównych pętli, które będę musiały po sobie następować w obrębie kodu danego typu warstwy, a ilość wymiarów każdej z macierzy oznacza ilość zagnieżdżonych pętli w głównej pętli danej macierzy.
W faktycznej części dekwantyzacji wag, wewnątrz ciała wspomnianych pętli, korzystając z drugiego, wewnętrznego wskaźnika pozycji, wydzielany jest każdy blok wag wraz z jego współczynnikiem skali. W przypadku kwantyzacji 4-bitowej występuje krok wydobycia obu liczb 4-bitowych z odczytanej liczby 8-bitowej. Następnie wszystkie wagi bloku są mnożone przez skalę i cofane jest ich przesunięcie z kroku kwantyzacji. Z takich zdekwantyzowanych wag jest tworzona macierz o odpowiednich wymiarach, która jest przypisywana do odpowiedniej macierzy warstwy modelu.
Po zakończeniu dekwantyzacji, funkcja zwraca gotowy model z przypisanymi wagami, który jest gotowy do użycia.

\chapter{Podsumowanie}
W pracy przedstawiono przegląd metod kompresji wag modeli sztucznej inteligencji ze szczególnym uwzględnieniem technik kwantyzacji oraz zaprezentowano autorskie narzędzie "Uni-Quant" do szybkiej kwantyzacji i dekwantyzacji modeli TensorFlow.

Główne cele pracy obejmowały: analizę podstawowych pojęć związanych z reprezentacją liczb i kwantyzacją, omówienie współczesnych rozwiązań stosowanych w praktyce inżynierskiej oraz opracowanie praktycznej biblioteki umożliwiającej przeprowadzenie kwantyzacji i odzyskanie wag w formacie użytecznym do wnioskowania.

W części teoretycznej uporządkowano pojęcia tensora, parametrów modelu, formatów liczbowych oraz zagadnienia związane z kalibracją i metrykami oceny jakości modeli. Omówiono również wpływ wyboru reprezentacji na rozmiar modelu, przepustowość pamięciową i czas wykonywania obliczeń.

Rozdział dotyczący popularnych rozwiązań zawierał charakterystykę ekosystemu współczesnych narzędzi do kwantyzacji. Przedstawiono podejścia takie jak BitsAndBytes (w tym format NF4 i LLM.int8()), algorytmy optymalizujące błąd rekonstrukcji wyjść (GPTQ), podejścia aktywacyjne uwzględniające ważność kanałów (AWQ) oraz nowsze metody typu TurboQuant. Omówiono także techniki redukcji metadanych, takie jak podwójna kwantyzacja, oraz formaty plików i praktyczne ograniczenia wdrożeń produkcyjnych.

W części implementacyjnej zaprezentowano bibliotekę Uni-Quant, która realizuje post-training quantization w wariantach 4- i 8-bitowych, opierając się na grupowaniu wag w bloki i zapisie współczynnika skalującego dla każdego bloku. Wprowadzono mechanizmy:
\begin{itemize}
    \item Blokowego skalowania wag oraz zapisu skal,\\
    \item Pakowania dwóch wartości 4-bitowych w jeden bajt,\\
    \item Zapisu wag i metadanych w postaci binarnej oraz JSON-owych manifestów scalonych w archiwum z rozszerzeniem .uniq,\\
    \item Zachowania fragmentów o rozmiarze mniejszym niż blok w postaci float32, aby zminimalizować degradację jakości.
\end{itemize}

Dodatkowo opracowano alternatywną implementację z wykorzystaniem CUDA, przenosząc krytyczne obliczenia do kerneli uruchamianych na GPU w celu przyspieszenia procesu kwantyzacji i dekwantyzacji. Rozwiązanie to umożliwia znaczące skrócenie czasu przetwarzania dużych modeli i stanowi praktyczny wybór tam, gdzie dostępne są akceleratory GPU.

Testy przeprowadzone w pracy wykazały oczekiwane zależności: zmniejszenie liczby bitów prowadzi do istotnego spadku rozmiaru pliku modelu, natomiast wpływ na jakość zależy od strategii grupowania wag (rozmiar bloku) oraz szczegółów kalibracji. Mniejszy rozmiar bloku zwykle zwiększa liczbę współczynników skalowania, co może prowadzić do większego rozmiaru pliku przy bardzo drobnej granularności, ale pozwala lepiej dopasować skalę lokalną i ograniczyć stratę jakości. Zachowanie niektórych małych fragmentów wag w formacie FP32 okazało się praktycznym kompromisem redukującym negatywny wpływ kwantyzacji na krytyczne parametry modelu.

Własny wkład praktyczny można podsumować jako:
\begin{itemize}
    \item Zaprojektowanie i zaimplementowanie biblioteki Uni-Quant do szybkiej kwantyzacji i dekwantyzacji modeli TensorFlow wraz z formatem archiwum .uniq,\\
    \item Przeprowadzenie testów ilustrujących kompromisy między rozmiarem pliku a zachowaniem jakości oraz wpływu rozmiaru bloku i szerokości bitowej,\\
    \item Integracja przyspieszenia obliczeń opartego na kernelach CUDA oraz udokumentowanie wpływu akceleracji na czasy przetwarzania.
\end{itemize}

Ograniczenia pracy i kierunki dalszych badań: praca koncentruje się na podejściach post-training quantization i prostych strategiach blokowego skalowania. Do istotnych rozszerzeń należą integracja zaawansowanych algorytmów PTQ (np. GPTQ i AWQ), implementacja per-channel scaling i technik kalibracji opartych na reprezentatywnych zbiorach danych. Dalsze badania mogą również objąć automatyczne dobieranie parametrów (np. rozmiar bloku) przy uwzględnieniu kosztu pamięciowego i metryk jakości oraz bardziej rozbudowane testy na modelach produkcyjnych o bardzo dużej liczbie parametrów. Praktycznym kierunkiem rozwoju jest bezpośrednia integracja formatu .uniq z ekosystemem TensorFlow, umożliwiająca wykonywanie modelu bezpośrednio z archiwum .uniq przy użyciu mechanizmu on‑the‑fly (lub częściowej, leniwej) dekwantyzacji bloków podczas inferencji, co poza obecnym zmniejszaniem rozmiaru na dysku, obniży wymagania sprzętowe i pozwoli uruchamiać duże modele w środowiskach o ograniczonych zasobach.

Podsumowując, praca ukazuje, że kwantyzacja jest efektywnym mechanizmem zmniejszania kosztów przechowywania i wykonywania modeli SI, lecz wymaga starannego doboru strategii i parametrów, aby zminimalizować spadek jakości. Uni-Quant dostarcza praktycznego narzędzia do eksperymentów i stanowi solidną bazę do dalszego rozwoju metod kompresji.

\nocite{*}
\printbibliography[heading=bibintoc,title={Bibliografia}]
\listoftables
\listoffigures
\end{document}